%% Согласно ГОСТ Р 7.0.11-2011:
%% 5.3.3 В заключении диссертации излагают итоги выполненного исследования, рекомендации, перспективы дальнейшей разработки темы.
%% 9.2.3 В заключении автореферата диссертации излагают итоги данного исследования, рекомендации и перспективы дальнейшей разработки темы.

\begin{enumerate}
  \item Установлено, что хотя прямые расчёты плотности и коэффициента самодиффузии демонстрируют расхождение с экспериментальными данными, последующая коррекция расчётных плотностей позволяет добиться более близкого к референтным значениям согласия по коэффициенту диффузии.
  \item Воспроизведено значение парахора чистного этана в диапазоне давлений 4–40 атм, что подтверждает возможность использования метода молекулярной динамики для прогнозирования экспериментальных значений поверхностного натяжения.
  \item Проведена верификация формулы Гиббса для оценки адсорбции на плоской поверхности в двухфазной двухкомпонентной системе. Формула, использующая только экспериментально измеряемые величины без приближений идеального газа и идеального раствора, успешно проверена на системе «метан–этан».
  \item Разработана и апробирована методика расчёта коэффициента термической аккомодации в рамках метода молекулярной динамики. Получены зависимости коэффициента от температуры и размера кластера для наночастиц железа, взаимодействующих с атомами аргона.
  \item Показано, что количество энергии, уносимой от кластера налетающим атомом, зависит от температуры кластера слабее, чем по линейному закону. Установлено, что данная нелинейность обусловлена уменьшением времени взаимодействия атома с кластером при росте температуры.
\end{enumerate}
